\newcommand{\figuraVarillaArticuladaPlomo}{%
\begin{figure}[H]
    \centering
    \begin{tikzpicture}[>=Latex,font=\small,line cap=round,line join=round]
        \fill[blue!12] (0.7,0.55) rectangle (11.4,3.75);
        \draw[very thick,blue!65!black] (0.7,3.75)--(11.4,3.75);
        \draw[very thick] (0.7,0.55)--(11.4,0.55);

        \coordinate (B) at (8.8,3.75);
        \coordinate (P) at (3.1,0.46);
        \draw[very thick,brown!80!black,line width=7pt] (B)--(P);
        \fill[white] (B) circle (0.12);
        \draw[very thick] (B) circle (0.12);
        \node[anchor=south west] at (B) {$B$};
        \draw[very thick] (8.8,3.75)--(9.65,4.45)--(9.65,3.05)--cycle;

        \fill[gray!65] (3.1,0.78) rectangle (3.75,1.42);
        \draw[very thick] (3.1,0.78) rectangle (3.75,1.42);
        \node[white,font=\scriptsize,align=center] at (3.425,1.10)
            {$2\,\mathrm{kg}$\\Pb};

        \draw[->,ultra thick,red!75!black] (5.55,2.25)--(5.55,3.45)
            node[above] {$\vec F_B$};
        \draw[->,ultra thick,orange!85!black] (6.35,2.38)--(6.35,1.20)
            node[below] {$\vec W$};
        \draw[->,thick,purple!75!black] (3.43,0.78)--(3.43,0.18)
            node[below] {$\vec W_{\mathrm{Pb}}$};

        \draw[densely dashed] (6.7,3.75)--(8.8,3.75);
        \draw[->,thick] (8.0,3.75)
            arc[start angle=180,end angle=210,radius=0.8];
        \node at (7.65,3.35) {$30^\circ$};
        \node[blue!65!black] at (10.55,2.9) {agua};
    \end{tikzpicture}
    \caption{Varilla uniforme articulada en $B$, sobre la línea de flotación,
    con una masa de plomo unida al extremo sumergido. La posición mostrada
    corresponde a $\theta=30^\circ$.}
    \label{fig:varilla-articulada-plomo}
\end{figure}%
}